\documentclass{situationist-preprint}

\title{Kane's Businesswoman Reconsidered}
\msid{ACT-15-b}
\msfamily{ACT-15}
\msdiscipline{Action, Agency and Free Will}
\mstype{Companion manuscript}
\mskeywords{Kane; luck; libertarianism; dual efforts; self-forming choice; contrastive explanation}

\begin{document}
\maketitle

\begin{abstract}
Kane's businesswoman, torn between helping the victim of an assault and hurrying to a meeting, is the standard case of a self-forming choice, and the standard target of the luck objection. We apply the result of ACT-15-a to the case. She is responsible for whichever choice she makes, because she tries to make each of them and each effort issues from a mechanism of her own. Two elements of Kane's presentation are not needed: the suggestion that she wills both outcomes as though in a single plan, and the neural model of her efforts. Without them the case is sound, and the luck objection fails because its demand for a contrastive explanation is not a condition of responsibility.
\end{abstract}

\section{The case and the objection}

\citet{Kane1996,Kane1999} describes a businesswoman on her way to an important meeting who sees an assault in an alley. Her moral convictions tell her to stop and help; her ambition tells her to go on. She struggles, and the outcome of the struggle is undetermined: with the same past and the same laws, she might have stopped or gone on. Kane holds that whichever she does, she is responsible, because she was trying to do each, and succeeded in the one she did through her own effort.

The luck objection, pressed by \citet{Mele2006}, asks what accounts for the difference between the actual world, in which she stops, and a world with the same past in which she goes on. Nothing about her accounts for it; so, the objection concludes, the difference is a matter of luck, and luck in what she does undermines her responsibility for it.

\section{The case reconsidered}

The case satisfies Definition~2.1 of ACT-15-a: the choice is between competing outcomes, one of which will occur, and which occurs is not determined by the past and the laws. We consider the two assumptions of Theorem~3.1 of ACT-15-a in turn.

\begin{proposition}\label{p:dual}
The businesswoman satisfies Assumption~2.3 of ACT-15-a: she tries to bring about each of the competing outcomes.
\end{proposition}

\begin{proof}
She is making an effort to overcome her ambition, which she believes is a way of stopping to help, and her effort is guided by stopping; she is also making an effort to resist the pull of her conscience, which she believes is a way of keeping her appointment, guided by going on. Each effort is a trying in the sense of Definition~2.2 of ACT-11-a. That each is opposed by the other does not prevent either from being a trying, as the effort of a player who aims at two targets at once is not undone by the other aim (Proposition~3.3 of ACT-15-a).
\end{proof}

Kane says that she wills each outcome. If her two volitions had to agglomerate into a single plan, she would be irrational, since no consistent plan includes both outcomes. By ACT-15-a this is not needed: the requirement that intentions be consistent with belief applies to each singly, and agglomeration is required only of intentions in one plan. Her two efforts belong to two plans, each of which she may abandon when the other succeeds.

\begin{proposition}\label{p:control}
The businesswoman satisfies Assumption~2.2 of ACT-15-a, and so, by Theorem~3.1 of ACT-15-a, she is responsible for whichever choice she makes.
\end{proposition}

\begin{proof}
Each of her efforts issues from her own reasons, her moral convictions on one side and her ambition on the other; each would vary with the considerations that support it, and neither was installed by another agent in a way that bypasses her reasoning. Her efforts therefore issue from a reasons-responsive mechanism of her own, which is Assumption~2.2. Together with Proposition~\ref{p:dual}, Theorem~3.1 of ACT-15-a applies.
\end{proof}

\section{The objection answered}

The luck objection asks for a contrastive explanation: what accounts for her stopping rather than going on. By Principle~P1 of ACT-15-a, responsibility for an outcome requires that it result from an effort the agent makes to bring it about, issuing from a mechanism of her own, even if the effort might have failed. It does not require that something about her explain why that effort succeeded rather than the other. The archer who aims well and hits is responsible for the hit, though given indeterminism in her release she might have missed. In the businesswoman's case there are two archers, both of them her, and whichever hits, it is her effort that hits.

Kane himself grants that no contrastive explanation of her choice is available, and rests her responsibility on her plural voluntary control, supported by a model of her two efforts as parallel neural processes. We agree that the contrastive demand need not be met, and Principle~P1 gives the reason without the neural model: the demand is not a condition of responsibility. The second element of Kane's presentation that is not needed is therefore the model, which describes how dual efforts might be realized but adds nothing to why they suffice.

The dual effort also marks the limit of the result. Had she tried only to reach the meeting, making no effort to stop, and had an undetermined fluctuation nonetheless halted her at the alley, the stopping would not have resulted from any effort of hers, Principle~P1 of ACT-15-a would not apply to it, and she would not be responsible for it.

\section{Conclusion}

The businesswoman tries to make each of the competing choices, each effort is her own, and she is responsible for whichever she makes. The case needs neither a single plan that includes both outcomes nor a neural model of her efforts, and no contrastive explanation.

\begin{thebibliography}{9}\small
\bibitem[Kane(1996)]{Kane1996} Kane, R. (1996). \emph{The Significance of Free Will}. Oxford University Press.
\bibitem[Kane(1999)]{Kane1999} Kane, R. (1999). Responsibility, luck, and chance: reflections on free will and indeterminism. \emph{The Journal of Philosophy}, 96(5), 217--240.
\bibitem[Mele(2006)]{Mele2006} Mele, A.~R. (2006). \emph{Free Will and Luck}. Oxford University Press.
\end{thebibliography}

\end{document}
